% =============================================================================
% VERSIÓN RBF/BJP – Doble columna (revtex4-2, estilo APS/PRE)
% Compilar: pdflatex borrador_rbf.tex  (dos pasadas)
% =============================================================================

\documentclass[aps,pre,twocolumn,reprint,superscriptaddress]{revtex4-2}

\usepackage[utf8]{inputenc}
\usepackage[T1]{fontenc}
\usepackage[spanish,es-nodecimaldot]{babel}
\usepackage{amsmath,amssymb}
\usepackage{graphicx}
\usepackage{siunitx}
\usepackage{booktabs}

% --- Macros físicas ----------------------------------------------------------
\newcommand{\cor}{e}
\newcommand{\coR}[1]{e_{#1}}
\newcommand{\ain}[1]{\alpha_{\mathrm{int},#1}}
\newcommand{\asu}[1]{\alpha_{\mathrm{sup},#1}}
\newcommand{\Ecm}[1]{E_{\mathrm{cm},#1}}
\newcommand{\Eint}[1]{E_{\mathrm{int},#1}}
\newcommand{\Tint}{T_{\mathrm{int}}}
\newcommand{\Tstar}{T^*}
\newcommand{\We}{\mathrm{We}}
\newcommand{\eps}{\varepsilon}
\newcommand{\epsps}{\varepsilon_{ps}}
\newcommand{\gstar}{g^*}
\newcommand{\lammps}{\textsc{Lammps}}

% =============================================================================
\begin{document}

\title{Partición de energía en el impacto de nanocúmulos
Lennard-Jones sobre superficies FCC: cuatro canales
y escalado con el tamaño del proyectil}

\author{Felipe González}
\email{fegonzalezw@gmail.com}
\affiliation{Pontificia Universidad Católica de Valparaíso (PUCV), Chile}

\author{Gonzalo [Apellido]}
\affiliation{Pontificia Universidad Católica de Valparaíso (PUCV), Chile}

\begin{abstract}
Estudiamos mediante dinámica molecular (\lammps{}, unidades LJ) la partición de
energía en el primer rebote de un nanocúmulo de Lennard-Jones sobre una
superficie FCC cristalina bajo campo gravitatorio, reproduciendo y extendiendo
el trabajo de Gutiérrez et al.\ (2012). Para un cúmulo de $N\approx83$ átomos
con interacción bola-superficie puramente repulsiva ($c_{ps}=0$, $T=0$),
identificamos cuatro canales energéticos distintos cuya suma es exactamente
igual a la energía de impacto $E_0$:
\[
    \alpha_{\mathrm{cm}} + \alpha_{T_B} + \alpha_{\Delta V_B}
    + \alpha_{\mathrm{surf}} = 1,
\]
donde $\alpha_{\mathrm{cm}} = e_1^2$ es la fracción recuperada como energía
traslacional del CM, $\alpha_{T_B}$ la KE vibracional interna de la bola,
$\alpha_{\Delta V_B}$ la variación neta de la PE intra-bola (expansión
anarmónica post-impacto), y $\alpha_{\mathrm{surf}}$ la excitación directa de
fonones de la superficie. El resultado central es que, para velocidades de
impacto $v_{\mathrm{imp}} \geq 0.9\,\sigma/\tau$, la expansión anarmónica
supera ampliamente a la excitación superficial directa:
$\alpha_{\Delta V_B} \approx 21$--$36\%$ frente a
$\alpha_{\mathrm{surf}} \approx 6$--$15\%$, revirtiendo la interpretación
cuantitativa de Gutiérrez et al., donde ambas contribuciones quedaban mezcladas
en un único canal ``superficial'' de 30--50\%.
La fracción cinética interna $\alpha_{T_B} \approx 0.22$ permanece constante
en ese rango de velocidades, consistente con el resultado original.

Extendemos el estudio a proyectiles de $N \in \{1, 2, 7, 13, 55, 83\}$ átomos.
Identificando la fracción interna total $\alpha_{\mathrm{bola}} = \alpha_{T_B}
+ \alpha_{\Delta V_B}$ con la ley de equipartición parcial de Gutiérrez,
$\alpha_{\mathrm{int}}(N) = (3N-3)/(3N-3+3N_c)$, se obtiene un número efectivo
de osciladores superficiales $N_c \approx 84 \approx N_{\mathrm{bola}}$, en
lugar de $N_c = 304$ que resulta de usar $\alpha_{T_B}$ solo. La interpretación
es que cada átomo de la bola se acopla efectivamente con un oscilador de la
superficie, lo que predice el límite asintótico
$\alpha_{\mathrm{bola}} \to \frac{1}{2}$ para $N \to \infty$.
Adicionalmente, la fracción superficial directa $\alpha_{\mathrm{surf}}(N)$
presenta un máximo no monótono en $N=13$ ($\approx 32\%$), correspondiente
al icosaedro de Mackay como tamaño de acoplamiento óptimo.

\medskip
\noindent\textit{Palabras clave:} dinámica molecular, nanocúmulos LJ,
coeficiente de restitución, partición energética, expansión anarmónica,
fonones de superficie.
\end{abstract}

\maketitle

% =============================================================================
\section{Introducción}
\label{sec:intro}
% =============================================================================

La pelota que rebota es uno de los arquetipos más estudiados de la física
disipativa. En el límite macroscópico y elástico, el coeficiente de restitución
$\cor = |v_f|/|v_i|$ es una propiedad del material que no cambia de un rebote
al siguiente: la altura disminuye como $h_n = h_0 \cor^{2n}$ y el cúmulo de
rebotes forma una serie geométrica convergente. Esta imagen simple descansa sobre
un supuesto crucial: que la energía disipada va \emph{íntegramente} a la
superficie (o al medio), sin que la pelota acumule calor interno entre rebote y
rebote.

A escala nanométrica este supuesto falla. Un cúmulo de $N$ átomos tiene $3N-6$
grados de libertad vibracionales internos que pueden absorber energía durante
cada impacto. Aunque el cúmulo comience en equilibrio térmico con la superficie
($\Tint = \Tstar$), cada rebote convierte parte de la energía mecánica
gravitatoria —que no forma parte del equilibrio térmico— en calor interno:
$\Tint(n) > \Tstar$ para todo $n \geq 1$ cuando $N \geq 2$. A medida que el
sobrecalentamiento $\Delta\Tint(n) = \Tint(n) - \Tstar$ crece, la densidad de
modos accesibles aumenta y la fracción de energía capturada internamente
$\ain{n}$ se incrementa en rebotes sucesivos. El resultado es una
\emph{cascada de disipación}: $\coR{1} > \coR{2} > \cdots$ con una tasa de
decaimiento que se acelera durante la secuencia. La fuente de energía de la
cascada es la gravedad, no la temperatura: el cúmulo extrae energía mecánica del
campo gravitatorio y la convierte en calor interno, bombeando el sistema fuera
del equilibrio en cada rebote.

Este fenómeno es cualitativamente distinto de la granular cooling de Haff
\cite{Haff1983}, donde un gas granular se enfría colectivamente por colisiones
inelásticas entre macropartículas sin estructura interna. Aquí la disipación tiene
dos canales: (i) transferencia a la red superficial (el mecanismo de Haff/Dahneke
\cite{Dahneke1975}) y (ii) almacenamiento en los modos internos del cúmulo, que
es reversible a escala de tiempo de la simulación pero irreversible en el sentido
estadístico una vez que $\Tint$ supera la escala de los fonones del cúmulo.

El trabajo más cercano al presente es el de Gutiérrez et al.~\cite{Gutierrez2012},
quienes simularon mediante dinámica molecular (código LPMD) el rebote repetido de
una bola FCC de argón de $N \approx 79$ átomos sobre una superficie del mismo
material, bajo un campo de fuerza constante equivalente a una gravedad reducida
$\gstar = 0.026\,\sigma/\tau^2$. Sus principales resultados fueron: (i) la
secuencia de alturas $h_n$ decrece más rápido que la serie geométrica predicha
por un COR constante; (ii) la energía interna de la bola —descompuesta en
energía cinética relativa al CM ($T_B$) y energía potencial intra-bola ($V_B$)—
aumenta en promedio con cada rebote; (iii) para alturas de caída suficientemente
grandes ($h_0 > 41\,\sigma$), la bola se funde durante el impacto. Sin embargo,
ese trabajo utilizó una interacción bola-superficie \emph{puramente repulsiva}
($c = 0$), por lo que la adhesión era imposible por construcción, y no estudió la
dependencia con el tamaño del cúmulo $N$ ni la influencia de la temperatura de
la superficie. Adicionalmente, su descomposición energética agrupa en un único
canal «superficial» la variación de la PE intra-bola y la excitación fonónica
directa de la red, lo que sobreestima esta última en un factor $\sim 3$--$5$.

Estudios de impacto único de cúmulos LJ sobre superficies FCC
\cite{Awasthi2007, Awasthi2010} han establecido la dependencia de $\cor(v_0, N,
\theta, \epsps)$ y la transición adhesión-rebote controlada por el número de Weber
$\We$. Sin embargo, la dinámica \emph{secuencial} bajo campo gravitatorio —la
pregunta de cuántos rebotes tarda un nanocúmulo en convertir su energía
mecánica en calor interno y eventualmente adherirse— no ha sido estudiada
sistemáticamente con interacción completa y temperatura finita.

Motivaciones concretas para este estudio incluyen:
\begin{itemize}
    \item \textbf{Deposición suave (\emph{soft landing})} de cúmulos sobre
    sustratos funcionales, donde el número de rebotes antes de la adhesión
    determina la morfología final del depósito \cite{Haberland1994}.
    \item \textbf{Fragmentación}: si $\Tint$ supera la temperatura de fusión
    del cúmulo antes de la adhesión, el cúmulo puede fragmentarse durante un
    rebote.
    \item \textbf{Fundamentos estadísticos}: la secuencia $\{\coR{n}\}$ es un
    experimento análogo a la relajación de un sistema de baja dimensionalidad
    ($3N-6$ modos) acoplado a un baño de calor macroscópico (la superficie).
\end{itemize}

En este trabajo reproducimos los experimentos de Gutiérrez et al.\ con
\lammps{} ($c_{ps}=0$, $T=0$, mismos parámetros) y extendemos su análisis en
dos direcciones. \textbf{(i)}~Descomponemos la energía del primer rebote en
\emph{cuatro} canales medidos de forma independiente, separando la variación
neta de la PE intra-bola ($\alpha_{\Delta V_B}$, expansión anarmónica
post-impacto) de la excitación directa de la superficie ($\alpha_{\mathrm{surf}}$),
con la identidad $\alpha_{\mathrm{cm}} + \alpha_{T_B} + \alpha_{\Delta V_B}
+ \alpha_{\mathrm{surf}} = 1$ exacta rebote a rebote.
\textbf{(ii)}~Extendemos el barrido a $N \in \{1, 2, 7, 13, 55, 83\}$ y
mostramos que la fórmula de equipartición de Gutiérrez describe los datos con
$N_c \approx N_{\mathrm{bola}}$ solo al usar la fracción interna total
$\alpha_{\mathrm{bola}} = \alpha_{T_B} + \alpha_{\Delta V_B}$, lo que
implica que cada átomo de la bola se acopla efectivamente con un oscilador
de la superficie y que $\alpha_{\mathrm{bola}} \to \tfrac{1}{2}$ para
$N \to \infty$.

% =============================================================================
\section{Modelo y definiciones}
\label{sec:modelo}
% =============================================================================

\subsection{Hamiltoniano del sistema y descomposición energética}
\label{sec:hamiltoniano}

El sistema consta de un cúmulo de $N$ átomos de tipo LJ (``proyectil'') y una
superficie cristalina FCC de $N_s$ átomos del mismo tipo (``superficie''). El
potencial de interacción entre cualquier par de átomos a distancia $r$ es
\begin{equation}
    U_{\mathrm{LJ}}(r) = 4\eps \left[ \left(\frac{\sigma}{r}\right)^{12}
    - c_{ij}\left(\frac{\sigma}{r}\right)^{6} \right],
    \quad r < r_c = 2.5\,\sigma,
    \label{eq:lj}
\end{equation}
donde $c_{ij} = 1$ para pares cúmulo-cúmulo y superficie-superficie, y
$c_{ij} = \sqrt{\epsps/\eps}$ para pares cúmulo-superficie, de modo que la
profundidad del pozo de adhesión es $\epsps$.

En unidades LJ ($\eps = \sigma = m = 1$), la energía total del cúmulo se
descompone exactamente como
\begin{equation}
    H_{\mathrm{cum}} = \underbrace{\frac{1}{2}M V_{\mathrm{cm}}^2 + MgZ_{\mathrm{cm}}}_{E_{\mathrm{mec}}}
    + \underbrace{\sum_{i=1}^{N}\frac{1}{2}m|\mathbf{v}_i - \mathbf{V}_{\mathrm{cm}}|^2
    + \sum_{i<j}U_{\mathrm{LJ}}(r_{ij})}_{E_{\mathrm{int}}},
    \label{eq:ham}
\end{equation}
donde $M = Nm$ es la masa total, $\mathbf{V}_{\mathrm{cm}}$ la velocidad del CM
y $g$ la aceleración gravitacional reducida. La energía mecánica del CM,
$E_{\mathrm{mec}}$, engloba la cinética traslacional y la energía potencial
gravitatoria. La energía interna $E_{\mathrm{int}}$ contiene las velocidades
\emph{relativas} al CM y las interacciones intra-cúmulo: es invariante bajo la
traslación y el estado de agitación interna del cúmulo.

Para $N=1$, $E_{\mathrm{int}} \equiv 0$ por definición: toda la energía del
cúmulo es mecánica. Este caso sirve como \textbf{referencia sin modos internos}.

\subsection{Protocolo de rebotes bajo gravedad}
\label{sec:protocolo}

\paragraph{Preparación.}
El cúmulo se construye a partir de la estructura de mínima energía global para
cada $N$ (icosaedros de Mackay para $N=13$ y $N=55$
\cite{Northby1987,Wales1997}) y se termaliza a $\Tint^{(0)} = \Tstar$ mediante
NVT de Nosé-Hoover durante $10^4\,\tau$ antes de iniciar la caída, en ausencia
de campo gravitatorio. Esta condición inicial coloca al cúmulo en equilibrio
térmico con la superficie: no hay diferencia de temperatura que pueda actuar
como fuente de la cascada; la única fuente de energía no-equilibrio es la
gravedad. La superficie FCC se equilibra de forma idéntica e independiente a
$\Tstar$ (ver Sección~\ref{sec:lammps}).

\paragraph{Condición inicial.}
El cúmulo se coloca a altura $Z_0$ sobre la superficie con velocidad CM nula
($\mathbf{V}_{\mathrm{cm}} = \mathbf{0}$) y temperatura interna $\Tint^{(0)} =
\Tstar$. El campo gravitatorio $\gstar$ lo acelera hacia la superficie; la
velocidad de impacto del primer rebote es
\begin{equation}
    v_{0,1} = \sqrt{2\gstar Z_0}.
\end{equation}
La energía mecánica inicial es $E_{\mathrm{mec}}^{(0)} = Mg\gstar Z_0$ y la
energía interna inicial es $\Eint{0} = \frac{3N-6}{2}k_B\Tstar$ (energía
térmica de equilibrio a $\Tstar$). El observable de interés no es $\Eint{}$
absoluta, sino el \emph{exceso}
\begin{equation}
    \Delta\Eint{n} \equiv \Eint{n} - \Eint{0}
    \label{eq:deltaE}
\end{equation}
generado por la cascada de rebotes.

\paragraph{Rebote $n$.}
Denominamos rebote $n$ al evento en que el CM del cúmulo pasa por un mínimo de
$Z_{\mathrm{cm}}(t)$. Justo antes del rebote $n$, la velocidad del CM es
$V_n = \sqrt{2\gstar h_{n-1}}$, donde $h_{n-1}$ es la altura máxima alcanzada
tras el rebote anterior ($h_0 = Z_0$). El coeficiente de restitución del rebote
$n$ se define como
\begin{equation}
    \coR{n} = \frac{|V_{\mathrm{cm},z}^{+}|}{|V_{\mathrm{cm},z}^{-}|},
    \label{eq:cor_n}
\end{equation}
donde los superíndices $\pm$ denotan el instante inmediatamente posterior y
anterior al rebote, detectado como el cruce $V_{\mathrm{cm},z} = 0$. La
altura máxima del rebote $n$ es
\begin{equation}
    h_n = \frac{(V_{\mathrm{cm},z}^{+})^2}{2\gstar} = h_{n-1}\,\coR{n}^2.
    \label{eq:altura}
\end{equation}

\subsection{Balance energético por rebote}
\label{sec:balance}

Durante el contacto del rebote $n$, la variación de energía mecánica del CM se
reparte entre dos canales: transferencia a la red superficial ($\Delta
E_{\mathrm{sup},n}$) y transferencia a los modos internos del cúmulo ($\Delta
E_{\mathrm{int},n}$). Definiendo las fracciones adimensionales
\begin{align}
    \ain{n} &= \frac{\Delta E_{\mathrm{int},n}}{E_{\mathrm{cm},n}^{-}},
    \label{eq:alpha_int} \\
    \asu{n} &= \frac{\Delta E_{\mathrm{sup},n}}{E_{\mathrm{cm},n}^{-}},
    \label{eq:alpha_sup}
\end{align}
donde $E_{\mathrm{cm},n}^{-} = \frac{1}{2}MV_n^2$ es la energía cinética del CM
antes del impacto, la conservación de energía total exige
\begin{equation}
    \coR{n}^2 + \ain{n} + \asu{n} = 1.
    \label{eq:balance}
\end{equation}
La identidad~\eqref{eq:balance} se verifica numéricamente para cada rebote con
residuo $< 10^{-6}$ (errores de redondeo en doble precisión), lo que valida la
descomposición energética.

El \textbf{exceso de energía interna} acumulado por la cascada tras $n$ rebotes es
\begin{equation}
    \Delta\Eint{n} = \sum_{k=1}^{n} \ain{k}\,E_{\mathrm{cm},k}^{-}
                   = \sum_{k=1}^{n} \ain{k}\,Mg\gstar h_{k-1}.
    \label{eq:Eint_acum}
\end{equation}
Siguiendo la notación de Gutiérrez et al.~\cite{Gutierrez2012}, la energía
interna se descompone en dos partes que oscilan entre sí durante el vuelo libre:
\begin{equation}
    E_{\mathrm{int}} = \underbrace{\sum_{i=1}^{N}\tfrac{1}{2}m|\mathbf{v}_i -
    \mathbf{V}_{\mathrm{cm}}|^2}_{T_B}
    + \underbrace{\sum_{i<j}U_{\mathrm{LJ}}(r_{ij})}_{V_B}.
    \label{eq:TB_VB}
\end{equation}
Durante el vuelo libre, $T_B$ y $V_B$ intercambian energía (oscilaciones internas
del cúmulo) mientras $E_{\mathrm{int}} = T_B + V_B$ permanece constante. Por
esto, el observable relevante para seguir la cascada es $E_{\mathrm{int}}$ total,
no $T_B$ solo. Cuando se promedia sobre un periodo de vibración interna completo,
$\langle T_B \rangle = \langle V_B \rangle = E_{\mathrm{int}}/2$ (equipartición),
lo que justifica la estimación de temperatura:
\begin{equation}
    \Delta\Tint(n) \equiv \Tint(n) - \Tstar
                   = \frac{2\,\Delta\Eint{n}}{(3N-6)\,k_B}
    \qquad (N \geq 3).
    \label{eq:Tint}
\end{equation}

\subsection{Número de Weber y criterio de adhesión}
\label{sec:weber}

La transición de rebote a adhesión ocurre cuando la energía cinética del CM es
insuficiente para superar el pozo de adhesión proyectil-superficie. El número
de Weber por rebote se define como
\begin{equation}
    \We_n = \frac{E_{\mathrm{cm},n}^{-}}{N\epsps}
           = \frac{Mg\gstar h_{n-1}}{N\epsps}.
    \label{eq:weber}
\end{equation}
Estudios de impacto único \cite{Awasthi2007} establecen un umbral
$\We_c \approx 1.5 \pm 0.3$. El cúmulo se adhiere al primer rebote $n$ tal que
$\We_n < \We_c$.

\subsection{Parámetros del barrido}
\label{sec:barrido}

\begin{table*}[ht]
\centering
\caption{Espacio paramétrico. \textbf{Bloque A}: reproducción de
\cite{Gutierrez2012} con parámetros originales. \textbf{Bloque B}: estudio
extendido. Cada condición del bloque B se repite con 3~réplicas de semillas
de velocidades distintas.}
\label{tab:barrido}
\begin{tabular}{llll}
\toprule
Parámetro & Símbolo & Bloque A (reproducción) & Bloque B (extensión) \\
\midrule
Tamaño del cúmulo    & $N$       & $\approx 79$ (FCC esférico)   & $1,\,2,\,7,\,13,\,55$ \\
Interacción bola-sup.& $c_{ps}$  & $0$ (repulsiva pura)          & $\sqrt{\epsps/\eps}$ \\
Gravedad reducida    & $\gstar$  & $0{,}026$                     & $0{,}005;\;0{,}01;\;0{,}02$ \\
Temperatura sup.     & $\Tstar$  & $0$ (NVE global)              & $0{,}1;\;0{,}3;\;0{,}5$ \\
Energía de adhesión  & $\epsps$  & --- (no aplica)               & $0{,}30;\;1{,}00$ \\
Altura inicial       & $Z_0$     & $\{11;\;18{,}6;\;28;\;41;\;58{,}6;\;79{,}5\}\,\sigma$
                                                                  & $20\,\sigma$ (fija) \\
$\Tint^{(0)}$        &           & $0$                           & $\Tstar$ \\
\bottomrule
\end{tabular}
\end{table*}

La altura inicial $Z_0 = 20\,\sigma$ con $\gstar = 0.01$ da una velocidad de
primer impacto $v_{0,1} = \sqrt{2 \times 0.01 \times 20} \approx 0.63\,\sigma/\tau$,
dentro del rango de impactos moderados. El valor de $\gstar$ no corresponde a
la gravedad real de argón ($g_{\mathrm{LJ}} \sim 10^{-8}\,\eps/\sigma m$), sino
que es un parámetro de control que establece la escala de energía gravitatoria
respecto a $\eps$.

% =============================================================================
\section{Detalles computacionales}
\label{sec:lammps}
% =============================================================================

Las simulaciones se realizan con \lammps~7~Feb~2024 \cite{Thompson2022} usando
el potencial LJ truncado con corrección de cola para las interacciones
superficie-superficie, y sin corrección para el par cúmulo-superficie (la
corrección asume homogeneidad del bulk que no se cumple en la interfaz).

\paragraph{Superficie.}
Red FCC de argón, $8 \times 8 \times 6$ celdas unitarias
($a_0 = 1.5496\,\sigma$), con condiciones de contorno periódicas en $xy$ y no
periódicas en $z$. Las 2 capas inferiores están fijas
(\texttt{setforce 0 0 0}). Las capas 3--5 se termostatizadas con NVT de
Nosé-Hoover a $\Tstar$ (amortiguamiento $\tau_{\mathrm{NVT}} = 0.5\,\tau$). La
capa superior (capa 6) evoluciona en NVE libre.

\paragraph{Cúmulo.}
El cúmulo se construye desde la estructura de mínima energía, se minimiza con
quench y se integra en NVE durante todo el vuelo e impacto (sin termostato):
la temperatura interna $\Tint$ es un \emph{resultado}, no una condición impuesta.

\paragraph{Campo gravitatorio.}
Se aplica el fix \texttt{gravity} de \lammps{} únicamente sobre los átomos del
cúmulo:
\begin{verbatim}
fix f_grav cluster gravity ${g} zdir -1.0
\end{verbatim}
La superficie no experimenta gravedad (sus capas inferiores están fijas o
termostatizadas).

\paragraph{Detección de rebotes.}
La posición $Z_{\mathrm{cm}}(t)$ se registra cada 10 pasos ($\Delta t = 10
\times 0.005\,\tau = 0.05\,\tau$). Los mínimos de $Z_{\mathrm{cm}}$ se
identifican en post-análisis mediante la condición $\dot{Z}_{\mathrm{cm}} = 0$
con cambio de signo. Para cada rebote se extrae $V_{\mathrm{cm},z}^{\pm}$ como
promedio sobre una ventana de $\pm 5$ pasos alrededor del mínimo, y $\Eint{n}$
como promedio de $E_{\mathrm{int}}(t)$ en la fase de vuelo libre posterior.

\paragraph{Criterio de adhesión.}
Un rebote se clasifica como adhesión si $Z_{\mathrm{cm}}(t) < z_{\mathrm{sup}} +
2\sigma$ durante más de $100\,\tau$ consecutivos sin rebotar (el cúmulo queda
atrapado en el pozo de adhesión).

% =============================================================================
\section{Resultados}
\label{sec:resultados}
% =============================================================================

\subsection{Reproducción de Gutiérrez et al.\ (2012): validación cruzada}
\label{sec:reproduccion}

Como primer paso, reproducimos los resultados del trabajo fundacional de
Gutiérrez et al.~\cite{Gutierrez2012} utilizando nuestro entorno
\lammps{}+Python en lugar del código LPMD original. Esto tiene dos objetivos:
(i) validar que nuestra cadena de análisis produce resultados físicamente
consistentes con los publicados, y (ii) establecer una línea base sobre la
cual medir el efecto de las extensiones que introducimos.

\paragraph{Parámetros de reproducción.}
Los parámetros se fijan para coincidir con los del capítulo~12 de
\cite{Gutierrez2012}: bola FCC de argón con $N \approx 79$ átomos ($r \approx
2.64\,\sigma$), construida como corte esférico de la red FCC; superficie FCC
de $\approx 1500$ átomos; interacción bola-superficie \textbf{puramente
repulsiva} ($c_{ps} = 0$ en la ec.~\eqref{eq:lj}); bola y superficie
equilibradas a $T = 0$ antes de la caída; gravedad $\gstar = 0.026\,\sigma/\tau^2$;
alturas iniciales $h_0 \in \{11.03,\;18.64,\;28.15,\;41.47,\;58.6,\;79.53\}\,\sigma$.
No se impone termostato en ninguna capa durante los rebotes (NVE global).

\paragraph{Resultados de validación.}
La Figura~\ref{fig:reprod_alturas} muestra las alturas máximas $h_n$ alcanzadas
por el CM de la bola tras cada rebote para las seis alturas iniciales.
Las curvas convergen al mismo valor asintótico $h_\infty \approx
1.0$--$1.5\,\sigma$ sobre el nivel de contacto efectivo $z_0 \approx
11.6\,\sigma$, equivalente a $Z_{\mathrm{cm}} \approx 12.7\,\sigma$ en
coordenadas absolutas de caja, en acuerdo razonable con el valor reportado
$\approx 11\,\sigma$ en \cite{Gutierrez2012}. Para $h_0 \geq 41\,\sigma$ la altura cae
más del 90\% en los primeros dos rebotes, consistente con la termalización interna
intensa descrita en el trabajo original.

\begin{figure*}[ht]
    \centering
    \includegraphics[width=0.9\textwidth]{../data/outputs/reprod_2012/fig_R1_alturas.pdf}
    \caption{Reproducción de la Fig.~7 de \cite{Gutierrez2012}.
    Alturas sucesivas $h_n$ del CM de la bola ($N\approx83$, $c_{ps}=0$, $T=0$) bajo
    gravedad $\gstar=0.026$, para seis alturas iniciales
    $h_0 \in \{8.5,\;16.2,\;25.7,\;38.9,\;56.4,\;77.2\}\,\sigma$.
    El acuerdo cuantitativo con el código independiente LPMD valida la cadena de análisis.}
    \label{fig:reprod_alturas}
\end{figure*}

La Figura~\ref{fig:reprod_energia} reproduce la Fig.~9 de \cite{Gutierrez2012}:
la evolución temporal de $K_B$ (cinética traslacional del CM), $T_B$ (cinética
interna), $V_B$ (potencial interna), $U_c$ (energía colisional bola-superficie)
y $U_F$ (energía gravitatoria). Se observan los mismos rasgos cualitativos:

\begin{enumerate}
    \item \textbf{$U_c$ domina durante el contacto}: la energía colisional
    almacena la mayor parte de la energía «disipada» durante el rebote; las
    demás contribuciones son menores.
    \item \textbf{$T_B$ oscila entre rebotes}: durante el vuelo libre, $T_B$
    y $V_B$ intercambian continuamente. Esto confirma que el observable
    físicamente significativo es $E_{\mathrm{int}} = T_B + V_B$, no $T_B$
    solo (ver Sección~\ref{sec:balance}).
    \item \textbf{$K_B$ decrece monotónamente rebote a rebote}: el pico de
    $K_B$ en cada rebote es menor que en el anterior, confiriendo a la
    secuencia de alturas su carácter no geométrico.
\end{enumerate}

\begin{figure*}[ht]
    \centering
    \includegraphics[width=0.9\textwidth]{../data/outputs/reprod_2012/fig_R2_energia_h0_28.15.pdf}
    \caption{Reproducción de la Fig.~9 de \cite{Gutierrez2012}. Evolución temporal de
    las energías del sistema para $h_0=28.15\,\sigma$ (caso representativo con tres
    rebotes visibles). Panel superior: $K_B$ (CM traslacional), $T_B$ (KE vibracional
    interna), $\Delta V_B$ (PE intra-bola) y $U_c$ (colisional). Panel inferior:
    $\Delta U_F$ (gravitatoria). La suma total se conserva a mejor que $10^{-4}\,\eps$
    por átomo en toda la simulación.}
    \label{fig:reprod_energia}
\end{figure*}

Adicionalmente, extraemos la secuencia $\{\coR{n}\}$ —
un observable que el trabajo original no reportó explícitamente—:
$\coR{1} = 0.66$, $\coR{2} = 0.56$, $\coR{3} = 0.80$
para $h_0 = 28.15\,\sigma$ (definido como $\coR{n} = \sqrt{h_n/h_{n-1}}$).
Este decaimiento supone la transición de una pelota
«casi elástica» ($\coR{1} \approx 0.91$ para $h_0 = 11\,\sigma$) a una que
pierde el 86\% de su energía en el segundo rebote
($\coR{2} \approx 0.37$ para $h_0 = 58.6\,\sigma$), en un rango de apenas
cuatro alturas de caída.
El decrecimiento $\coR{1} > \coR{2} > \cdots$ (visible en la
Fig.~\ref{fig:reprod_alturas}) es la firma de la cascada de disipación predicha
por Gutiérrez: en cada rebote la bola llega con mayor energía interna acumulada
$E_{\mathrm{int},n}$, lo que incrementa $\alpha_{\mathrm{int},n}$ en rebotes
sucesivos. Sin embargo, cuantificar la partición energética rebote a rebote
—es decir, determinar $\alpha_{T_B}(n)$, $\alpha_{\Delta V_B}(n)$ y
$\alpha_{\mathrm{surf}}(n)$ para $n \geq 2$ descontando la energía interna
preexistente— requiere integrar el formalismo de la Sección~\ref{sec:balance}
sobre la trayectoria completa con $E_{\mathrm{int},0} > 0$. El presente trabajo
se concentra en $n=1$ (bola fría, $E_{\mathrm{int},0} = 0$), que es el caso
más limpio y sirve de referencia energética para el estudio de la cascada completa.

\paragraph{Régimen de alta energía y termalización interna.}
Para $h_0 \geq 41.47\,\sigma$, la bola almacena una energía interna
$T_B \geq 28\,\eps$ tras el primer rebote, suficiente para fundir la
estructura FCC (consistente con la Fig.~8 de \cite{Gutierrez2012}).
La altura predita por equipartición es $h^* = T_B/(M_b g) \approx 13\,\sigma$,
pero convergir a este plateau requiere cientos de rebotes más allá de
los $400\,\tau$ simulados. Nuestra reproducción captura correctamente el mecanismo
de termalización interna y su umbral en $h_0$.

\medskip
Con la validación establecida, la Sección~\ref{sec:4canales} presenta la
descomposición del primer rebote en cuatro canales energéticos (resultado
nuevo respecto de \cite{Gutierrez2012}), y la Sección~\ref{sec:sweep_N} reporta
el escalado de la fracción interna con el tamaño $N$ del proyectil.

\subsection{Partición energética del primer rebote en cuatro canales}
\label{sec:4canales}

El análisis de Gutiérrez et al.~\cite{Gutierrez2012} descompone la energía del
primer impacto en dos canales: fracción interna
$\alpha_{T_B} = T_B^{(1)}/E_0$ (KE vibracional de la bola medida al
primer apogeo) y fracción superficial $1 - e_1^2 - \alpha_{T_B}$ (residuo).
Esta definición asigna implícitamente al canal superficial todo lo que no es
CM ni KE interna.

Disponiendo de mediciones directas de la energía de la superficie
($K_S + V_{\mathrm{surf}}$), identificamos que el residuo mezcla dos
contribuciones físicamente distintas:
\begin{enumerate}
    \item $\alpha_{\Delta V_B}$: variación neta de la energía potencial
    intra-bola entre el estado inicial (red FCC a $T=0$) y el primer apogeo.
    Mide la expansión anarmónica del cúmulo que quedó sobre-comprimido
    durante el impacto y se expande al rebotar.
    \item $\alpha_{\mathrm{surf}}$: excitación directa de la superficie,
    medida como $\Delta(K_S + V_{\mathrm{surf}})/E_0$.
\end{enumerate}
Con la normalización CM ($\alpha_{\mathrm{cm}} = e_1^2$), la identidad
energética es exacta rebote a rebote:
\begin{equation}
    \alpha_{\mathrm{cm}} + \alpha_{T_B} + \alpha_{\Delta V_B}
    + \alpha_{\mathrm{surf}} = 1.
    \label{eq:4canales}
\end{equation}

La Tabla~\ref{tab:particion4} y la Figura~\ref{fig:particion4} muestran
esta descomposición para $N\approx83$, $c_{ps}=0$, $T=0$, en función de la
velocidad de impacto $v_{\mathrm{imp}} = \sqrt{2gh_0^{\mathrm{ef}}}$.

\begin{table*}[ht]
\centering
\caption{Partición de la energía del primer rebote en cuatro canales para
$N\approx83$, $c_{ps}=0$, $T=0$. Suma = 1 por construcción.}
\label{tab:particion4}
\begin{tabular}{cccccccc}
\toprule
$h_0^{\mathrm{ef}}$ & $v_{\mathrm{imp}}$ & $e_1$ &
$\alpha_{\mathrm{cm}}$ & $\alpha_{T_B}$ & $\alpha_{\Delta V_B}$ &
$\alpha_{\mathrm{surf}}$ \\
($\sigma$) & ($\sigma/\tau$) & & & & & \\
\midrule
 8.53 & 0.666 & 0.906 & 0.822 & 0.068 & 0.046 & 0.065 \\
16.21 & 0.918 & 0.697 & 0.486 & 0.211 & 0.247 & 0.056 \\
25.73 & 1.157 & 0.661 & 0.437 & 0.178 & 0.289 & 0.096 \\
38.94 & 1.423 & 0.660 & 0.435 & 0.223 & 0.214 & 0.128 \\
56.41 & 1.713 & 0.618 & 0.382 & 0.219 & 0.300 & 0.099 \\
77.24 & 2.004 & 0.503 & 0.253 & 0.233 & 0.361 & 0.153 \\
\bottomrule
\end{tabular}
\end{table*}

El resultado principal es que para $v_{\mathrm{imp}} \geq 0.9\,\sigma/\tau$,
$\alpha_{\Delta V_B} \gg \alpha_{\mathrm{surf}}$: la expansión anarmónica de
la bola absorbe entre el 21\% y el 36\% de $E_0$, mientras que la excitación
directa de la superficie capta solo entre el 6\% y el 15\%.
Esto revierte la interpretación cuantitativa de \cite{Gutierrez2012}, donde
$\alpha_{\mathrm{surf}}$ se estimaba en 30--50\% al englobar implícitamente
$\alpha_{\Delta V_B}$ en el canal superficial.

La fracción cinética interna $\alpha_{T_B}$ es prácticamente constante para
$v_{\mathrm{imp}} \geq 0.9\,\sigma/\tau$ ($\alpha_{T_B} \approx 0.21$--$0.23$),
consistente con el hallazgo original de Gutiérrez. Sin embargo, la fracción
interna total
\begin{equation}
    \alpha_{\mathrm{bola}} \equiv \alpha_{T_B} + \alpha_{\Delta V_B}
\end{equation}
crece del 46\% al 59\% en ese rango, y para la velocidad máxima estudiada
($v_{\mathrm{imp}}=2.0\,\sigma/\tau$) supera a todos los demás canales combinados.

\begin{figure*}[ht]
    \centering
    \includegraphics[width=0.85\textwidth]{../data/outputs/reprod_2012/fig_A2_particion.pdf}
    \caption{Partición de la energía del primer rebote en cuatro canales para
    $N\approx83$, $c_{ps}=0$, $T=0$, en función de la altura inicial $h_0$.
    Barras azules: fracción CM ($\alpha_{\mathrm{cm}} = e_1^2$). Rojas:
    KE interna ($\alpha_{T_B}$). Naranjas: PE anarmónica ($\alpha_{\Delta V_B}$).
    Grises: excitación directa de la superficie ($\alpha_{\mathrm{surf}}$). Las
    etiquetas dentro de cada segmento muestran el valor porcentual.}
    \label{fig:particion4}
\end{figure*}

\subsection{Variación sistemática con \texorpdfstring{$N$}{N}: escalado de
\texorpdfstring{$\alpha_{\mathrm{bola}}$}{alpha\_bola} y calibración de
\texorpdfstring{$N_c$}{Nc}}
\label{sec:sweep_N}

Gutiérrez et al.~\cite{Gutierrez2012} propusieron que la fracción interna del
primer rebote sigue la ley de equipartición parcial entre los $3N-3$ modos de la
bola y los $3N_c$ modos efectivos de la superficie de contacto:
\begin{equation}
    \alpha_{\mathrm{int}}(N) = \frac{3N-3}{3N-3 + 3N_c}.
    \label{eq:gutierrez}
\end{equation}
El parámetro $N_c$ (número efectivo de osciladores de superficie en la huella)
debe calibrarse con datos de simulación.

Corremos las mismas seis alturas con proyectiles de $N \in \{1, 2, 7, 13, 55, 83\}$
en el régimen $c_{ps}=0$, $T=0$, y promediamos sobre las alturas con
$v_{\mathrm{imp}} > 0.9\,\sigma/\tau$ para evitar el régimen de baja energía
donde la partición es distinta. Los resultados se recogen en la
Tabla~\ref{tab:sweep_N}.

\begin{table*}[ht]
\centering
\caption{Fracciones de energía del primer rebote promediadas sobre
$v_{\mathrm{imp}}>0.9\,\sigma/\tau$, en función de $N$. Las dos últimas columnas
muestran la predicción de la ec.~\eqref{eq:gutierrez} para los dos valores de
$N_c$ calibrados.}
\label{tab:sweep_N}
\begin{tabular}{ccccccc}
\toprule
$N$ & $e_1$ & $\alpha_{T_B}$ & $\alpha_{\mathrm{bola}}$ & $\alpha_{\mathrm{surf}}$ &
$\hat\alpha(N_c{=}304)$ & $\hat\alpha(N_c{=}84)$ \\
\midrule
 1 & 0.934 & 0.000 & 0.000 & 0.126 & 0.000 & 0.000 \\
 2 & 0.916 & 0.002 & 0.002 & 0.156 & 0.003 & 0.012 \\
 7 & 0.879 & 0.023 & 0.052 & 0.267 & 0.019 & 0.067 \\
13 & 0.763 & 0.053 & 0.086 & 0.318 & 0.038 & 0.125 \\
55 & 0.585 & 0.190 & 0.410 & 0.231 & 0.151 & 0.392 \\
83 & 0.628 & 0.213 & 0.495 & 0.106 & 0.213 & 0.495 \\
\bottomrule
\end{tabular}
\end{table*}

\paragraph{Calibración de $N_c$.}
Usando $N=83$ como referencia e identificando
$\alpha_{\mathrm{int}}$ con $\alpha_{T_B} = 0.213$, la ec.~\eqref{eq:gutierrez}
da $N_c = 304$, un valor sin interpretación geométrica obvia. En cambio,
identificando $\alpha_{\mathrm{int}}$ con la fracción interna total
$\alpha_{\mathrm{bola}} = 0.495$, se obtiene
\begin{equation}
    N_c = (3N-3)\left(\frac{1}{\alpha_{\mathrm{bola}}} - 1\right)/3
    \approx 84 \approx N_{\mathrm{bola}}.
    \label{eq:Nc_calibrado}
\end{equation}
Este resultado sugiere que, en promedio, \emph{cada átomo de la bola se acopla
efectivamente con un oscilador de la superficie}. Para $N_c = N$ la
ec.~\eqref{eq:gutierrez} predice el límite asintótico
\begin{equation}
    \alpha_{\mathrm{bola}}(N) = \frac{N-1}{2N-1}
    \xrightarrow{N\to\infty} \frac{1}{2},
    \label{eq:limite_N}
\end{equation}
lo que equivale a equipartición entre un sistema de $N$ osciladores
(la bola) y uno de $N$ osciladores efectivos (la huella en la superficie).

\paragraph{Comportamiento no monótono de $\alpha_{\mathrm{surf}}(N)$.}
Un resultado sorprendente en la Tabla~\ref{tab:sweep_N} es que
$\alpha_{\mathrm{surf}}$ no es monótona: crece de $N=1$ a $N=13$ (máximo
$\approx 32\%$) y luego decrece para $N=55$ ($23\%$) y $N=83$ ($11\%$).
Para $N$ pequeño, la bola es demasiado ligera para excitar fonones
superficiales eficientemente. Para $N$ grande, los modos internos absorben la
energía antes de que el impulso se transfiera a la red. El cúmulo de $N=13$
(icosaedro de Mackay con $r \approx 1.5\,\sigma$) representa el tamaño de
acoplamiento óptimo bola-superficie bajo las condiciones del estudio.

La Figura~\ref{fig:B_consolid} resume los cuatro observables en escala
logarítmica. Los paneles (c) y (d) comparan la predicción de la
ec.~\eqref{eq:gutierrez} con $N_c=304$ y $N_c=84$: la calibración con
$\alpha_{\mathrm{bola}}$ describe cuantitativamente los datos en todo el
rango $N\in\{2,\ldots,83\}$.

\begin{figure*}[ht]
    \centering
    \includegraphics[width=0.9\textwidth]{../data/outputs/cluster_sweep/fig_B_consolidada.pdf}
    \caption{Barrido $N \in \{1,2,7,13,55,83\}$ con $c_{ps}=0$, $T=0$,
    promediado sobre $v_{\mathrm{imp}}>0.9\,\sigma/\tau$.
    (a)~$e_1(N)$: decrece con $N$.
    (b)~$\alpha_{\mathrm{surf}}(N)$: no monótona, máximo en $N=13$.
    (c)~$\alpha_{T_B}(N)$ con ajuste de la ec.~\eqref{eq:gutierrez} para
    $N_c=304$ (línea punteada).
    (d)~$\alpha_{\mathrm{bola}}(N)$ con ajuste para $N_c=84\approx N_{\mathrm{bola}}$
    (línea sólida). El punto $N=7$ en (d) se indica con una flecha porque
    $E_0(N{=}7)\lesssim1\,\eps$, límite de relación señal-ruido de la
    medición de la energía de la superficie.}
    \label{fig:B_consolid}
\end{figure*}

% =============================================================================
\section{Discusión}
\label{sec:discusion}
% =============================================================================

\subsection{Origen físico de \texorpdfstring{$\alpha_{\Delta V_B}$}{alpha\_DVB}:
expansión anarmónica post-impacto}
\label{sec:anarmonico}

El canal $\alpha_{\Delta V_B}$ no tiene análogo en modelos de impacto de esferas
rígidas ni en la aproximación armónica. Su origen es la \emph{asimetría temporal}
del proceso de compresión-expansión del cúmulo LJ durante el contacto.

Durante la compresión (cúmulo decelerando hacia la superficie), el impulso de
la pared convierte KE del CM en energía potencial de la red intra-bola. A las
velocidades estudiadas ($v_{\mathrm{imp}} \geq 0.9\,\sigma/\tau$), la compresión
supera el rango harmónico y la bola queda en un estado de PE mayor que su
referencia de equilibrio. Al rebotar, la KE del CM se restaura, pero la relajación
de la PE interna hacia el equilibrio tarda varias vibraciones internas
($\sim 1$--$5\,\tau$), tiempo durante el cual el cúmulo ya se aleja de la
superficie. En el primer apogeo, $V_B > V_B^{(0)}$ y la diferencia
$\alpha_{\Delta V_B} = (V_B^{(1)} - V_B^{(0)})/E_0$ registra esta energía
potencial que la bola «se lleva» sin haber convertido aún en KE interna.

Este mecanismo tiene dos consecuencias:
\begin{enumerate}
    \item La descomposición de Gutiérrez —que toma $T_B^{(1)}$ al primer apogeo
    como medida de toda la energía interna— subestima el calor depositado en la
    bola. A velocidades altas, $\alpha_{\Delta V_B} > \alpha_{T_B}$: la PE supera
    a la KE interna en el apogeo.
    \item $\alpha_{\mathrm{bola}} = \alpha_{T_B} + \alpha_{\Delta V_B}$ es la
    cantidad conservada que entra al balance de canales; $\alpha_{T_B}$ solo es
    la fracción que ya se ha termalizado en el momento de la medición. En
    rebotes posteriores, $\alpha_{\Delta V_B}$ se redistribuirá entre los modos
    internos de la bola y eventualmente fluirá parcialmente de vuelta a la
    superficie.
\end{enumerate}

\subsection{Conexión con la ley de Haff}
\label{sec:haff}

La ley de Haff~\cite{Haff1983} describe el enfriamiento de un gas granular:
$T_{\mathrm{gran}}(t) \propto (1 + t/t_0)^{-2}$. En nuestro sistema el
análogo sería el decaimiento de $\Ecm{n}$ con $n$: para el caso $N=1$ con $\cor$
constante, $\Ecm{n} = \Ecm{0}\cor^{2n}$ es más bien una ley exponencial que una
ley de potencias.

La diferencia fundamental con Haff es que en un gas granular \emph{todas} las
partículas disipan entre sí y no hay reservorio de calor; la temperatura global
cae. En nuestro sistema, el calor no se pierde sino que se transfiere al interior
del cúmulo ($\Tint$ sube) y a la superficie termostatizada ($\Tstar$ fijo). La
irreversibilidad no viene de la entropía del gas granular sino de la diferencia
de temperatura $\Tint - \Tstar$ que impulsa el flujo de calor de vuelta a la
superficie en cada impacto.

\subsection{Números mágicos e \texorpdfstring{$\alpha_{\mathrm{surf}}$}{alpha\_surf}
no monótona}
\label{sec:magicos}

Los cúmulos $N=13$ y $N=55$ son icosaedros de Mackay con temperatura de fusión
local superior a la de sus tamaños vecinos. En las fracciones
$\alpha_{T_B}(N)$ y $\alpha_{\mathrm{bola}}(N)$ (paneles (c) y (d) de la
Fig.~\ref{fig:B_consolid}), ambos tamaños siguen la tendencia suave predicha
por la ec.~\eqref{eq:gutierrez} con $N_c=84$ sin desviaciones estadísticamente
significativas. La mayor estabilidad estructural de los icosaedros no reduce la
captura de energía interna en el primer impacto a las velocidades estudiadas.

La excepción es $\alpha_{\mathrm{surf}}$, que sí muestra una anomalía en $N=13$:
el máximo del panel (b) de la Fig.~\ref{fig:B_consolid} ($\approx 32\%$) supera
en un factor $\sim 3$ el valor para $N=83$ ($11\%$). La interpretación es
geométrica: el icosaedro de $N=13$ tiene un radio de contacto
$r_{13} \approx 1.5\,\sigma$ comparable al rango de los modos fonónicos de
superficie más acoplados ($\lambda \sim 2$--$3\,\sigma$), maximizando la
transferencia directa. Para $N \gg 13$, la zona de contacto crece pero los
modos internos del cúmulo absorben el exceso antes de que llegue a la red.

\subsection{Interpretación de \texorpdfstring{$N_c \approx N_{\mathrm{bola}}$}{Nc aprox Nbola}}
\label{sec:Nc_interpretacion}

El resultado $N_c = 84 \approx N_{\mathrm{bola}}$ tiene una lectura directa en
términos de la fórmula de equipartición ec.~\eqref{eq:gutierrez}: la energía
de impacto se distribuye entre $3N_{\mathrm{bola}}-3$ osciladores internos y
$3N_c \approx 3N_{\mathrm{bola}}$ osciladores externos (en la superficie),
es decir, la huella activa del impacto contiene aproximadamente el mismo número
de osciladores que la propia bola. Esta escala es consistente con la observación
de que la zona de contacto bola-superficie involucra del orden de $N^{2/3}$
átomos superficiales, pero que la profundidad de penetración activa varios planos
de la red, resultando en un número efectivo $\sim N$.

Para $N_c = N$, el límite asintótico $\alpha_{\mathrm{bola}} \to \tfrac{1}{2}$
es el resultado de equipartición entre dos sistemas de igual tamaño: es la
misma física que en la colisión elástica de dos masas iguales
($e=1 \Rightarrow$ toda la energía al impactado), con la diferencia de que aquí
la transferencia es irreversible porque los modos se excitan en formas normales
distintas a las del CM.

% =============================================================================
\section{Conclusiones}
\label{sec:conclusiones}
% =============================================================================

Hemos reproducido los resultados de Gutiérrez et al.~\cite{Gutierrez2012}
con \lammps{} ($c_{ps}=0$, $T=0$, $N\approx83$, $\gstar=0.026$) y extendido
su análisis energético en dos direcciones. Las conclusiones principales son:

\begin{enumerate}

    \item \textbf{El impacto involucra cuatro canales energéticos distintos.}
    La identidad $\alpha_{\mathrm{cm}} + \alpha_{T_B} + \alpha_{\Delta V_B}
    + \alpha_{\mathrm{surf}} = 1$ es exacta rebote a rebote y está respaldada
    por medición directa de la energía de la superficie. La descomposición en
    dos canales de Gutiérrez mezcla $\alpha_{\Delta V_B}$ con
    $\alpha_{\mathrm{surf}}$, sobreestimando la excitación superficial en un
    factor 3--5.

    \item \textbf{La expansión anarmónica domina sobre la excitación superficial
    directa.} Para $v_{\mathrm{imp}} \geq 0.9\,\sigma/\tau$,
    $\alpha_{\Delta V_B} \approx 21$--$36\%$ supera a
    $\alpha_{\mathrm{surf}} \approx 6$--$15\%$. Este canal refleja que el
    cúmulo LJ abandona el contacto en un estado de PE intra-bola mayor que
    el de equilibrio, y esa energía permanece almacenada al primer apogeo.
    $\alpha_{T_B} \approx 0.22$ es prácticamente constante en ese rango,
    consistente con Gutiérrez.

    \item \textbf{El número efectivo de osciladores superficiales es
    $N_c \approx N_{\mathrm{bola}}$.} Al calibrar la fórmula de Gutiérrez con
    la fracción interna total $\alpha_{\mathrm{bola}} = \alpha_{T_B} +
    \alpha_{\Delta V_B}$, se obtiene $N_c = 84 \approx N_{\mathrm{bola}}$,
    frente a $N_c = 304$ al usar $\alpha_{T_B}$ solo. El resultado implica
    que cada átomo de la bola se acopla efectivamente con un oscilador de la
    superficie, y que para $N \to \infty$ se alcanza el límite de equipartición
    $\alpha_{\mathrm{bola}} \to \tfrac{1}{2}$.

    \item \textbf{$\alpha_{\mathrm{surf}}(N)$ tiene un máximo no monótono en
    $N=13$.} La fracción de excitación fonónica directa alcanza $\approx 32\%$
    para el icosaedro de Mackay ($N=13$) y cae a $\approx 11\%$ para $N=83$.
    Este máximo coincide con la escala geométrica de acoplamiento óptimo entre
    el radio de contacto de la bola y los modos fonónicos de superficie.

    \item \textbf{El proyectil monoatómico ($N=1$) es la única referencia
    clásica.} Con $\alpha_{\mathrm{bola}} = 0$ por construcción,
    $e_1 = 0.934$ es estrictamente independiente de $h_0$ en el rango
    repulsivo ($c_{ps}=0$) y sirve como línea base para cuantificar el efecto
    de los modos internos en tamaños mayores.

\end{enumerate}

Como trabajo futuro inmediato, la extensión a interacción LJ completa
($c_{ps}=1$) y temperatura inicial $\Tstar > 0$ permitirá estudiar la
\emph{cascada de rebotes sucesivos} hasta la adhesión: el número de rebotes
$n_{\mathrm{adh}}(N)$, el sobrecalentamiento de pico $\Delta\Tint^{\mathrm{pico}}$
y la influencia de $\Tstar$ en la tasa de la cascada. Los resultados del
presente trabajo —en particular la calibración $N_c \approx N_{\mathrm{bola}}$—
proporcionan la referencia energética en $T=0$ sobre la que se medirán esos
efectos térmicos.

\begin{acknowledgments}
Los autores agradecen al programa de Magíster en Simulación Computacional de la
PUCV por el apoyo académico. Las simulaciones se realizaron con
\lammps~7~Feb~2024 \cite{Thompson2022}. El código fuente está disponible en
\url{https://github.com/fwasaff/inelastic-rebound-md}.
\end{acknowledgments}

% =============================================================================
\begin{thebibliography}{99}

\bibitem{Gutierrez2012}
G. Gutiérrez, S. Davis, C. Loyola, J. Peralta, F. González, Y. Navarrete
y F. González-Wasaff,
\textit{Inelastic Collisions and Hypervelocity Impacts at Nanoscopic Level:
A Molecular Dynamics Study},
en: \textit{Molecular Dynamics -- Theoretical Developments and Applications
in Nanotechnology and Energy}, cap.~12, pp.~229--248, InTech (2012).

\bibitem{Haff1983}
P.~K. Haff,
\textit{Grain flow as a fluid-mechanical phenomenon},
J. Fluid Mech. \textbf{134}, 401 (1983).

\bibitem{Dahneke1975}
B. Dahneke,
\textit{Further measurements of the bouncing of small latex spheres},
J. Colloid Interface Sci. \textbf{51}, 58 (1975).

\bibitem{Awasthi2007}
A. Awasthi, S.~C. Hendy, P. Zoontjens y S.~A. Brown,
\textit{Molecular dynamics simulations of reflection and adhesion behavior in
Lennard-Jones cluster deposition},
Phys. Rev. B \textbf{76}, 115437 (2007).

\bibitem{Awasthi2010}
A. Awasthi, S.~A. Brown y S.~C. Hendy,
\textit{Oblique impacts of Lennard-Jones clusters},
Phys. Rev. B \textbf{82}, 073406 (2010).

\bibitem{Northby1987}
J.~A. Northby,
\textit{Structure and binding of Lennard-Jones clusters: $13 \le N \le 147$},
J. Chem. Phys. \textbf{87}, 6166 (1987).

\bibitem{Wales1997}
D.~J. Wales y J.~P.~K. Doye,
\textit{Global optimization by basin-hopping and the lowest energy structures
of Lennard-Jones clusters containing up to 110 atoms},
J. Phys. Chem. A \textbf{101}, 5111 (1997).

\bibitem{Haberland1994}
H. Haberland, M. Mall, M. Moseler, Y. Qiang, T. Reiners y Y. Thurner,
\textit{Filling of micron-sized contact holes with copper by energetic cluster
impact},
J. Vac. Sci. Technol. A \textbf{12}, 2925 (1994).

\bibitem{Thompson2022}
A.~P. Thompson \textit{et al.},
\textit{LAMMPS -- A flexible simulation tool for particle-based materials
modeling},
Comput. Phys. Commun. \textbf{271}, 108171 (2022).

\end{thebibliography}

\end{document}
